\section{CNN}

\begin{definition}[Integral Op]
    Ker \(G: \R^2 \to \R\), \(t_1, t_2 \in \R_\infty\), then
    \((T f)(u) = \int_{t_1}^{t_2} G(u, t) f(t) \, dt\)
\end{definition}

\begin{definition}[Fourier]
    \((\mathcal F f)(u) = \int_\R f(t) e^{-i 2 \pi u t} \, dt\) \ \textbf{Conv}:
    \((f * g)(u) = \int_{-\infty}^{\infty} g(u - t) f(t) \, dt\)
\end{definition}

\begin{definition}[Time-shift]
    \(f_\Delta(t) := f(t + \Delta)\), then \(f_\Delta * h = (f * h)_\Delta\).
\end{definition}

\begin{theorem}[Convolution]
    \(\mathcal F(f * h) = \mathcal F f \cdot \mathcal F h\)
    \adv{\(\OB(n \log n)\) for FFT}
    w/ inverse: \(f * h = \mathcal F^{-1}((\mathcal F f) \cdot (\mathcal F h))\).
     \disadv{\(m < \log n\): favor time/space instead of freq}
\end{theorem}

\begin{definition}[Linear shift-equivariant]
    \(\forall f, g, \alpha, \beta: T(\alpha f + \beta g) = \alpha Tf + \beta Tg\)
    and \\ \((T f_\Delta)(t) = (Tf)(t + \Delta)\). \adv{Can \textit{always} be written as convolution}.
\end{definition}

\begin{definition}[Discrete Correlation]
    \(I'(i, j) = \sum_{m=-k}^k \sum_{n=-k}^k K(m, n) I(i+m, j+n)\)
\end{definition}

\begin{definition}[Discrete Convolution]
    \(I'(i, j) = \sum_{m=-k}^k \sum_{n=-k}^k I(i -m, j - n) K(m, n)\)

    \begin{itemize*}
        \item Corr \(\iff K(m, n) = K(-m, -n)\)
        \item Commutative
        \item Toepliz MM \(\in \R^{(D-K+1) \times D}\)
    \end{itemize*}
\end{definition}

\begin{definition}[CNN-formulas]
    \(\bm{C}\)hannels, \(\bm{K}\)ernel size, \(m = \#\)Kernels, \(\bm{{D}}\)ill., \(\bm{P}\)ad, \(\bm{S}\)tride
    \begin{itemize*}
        \item Dim: \(f(i) = \lfloor\frac{i + 2P_i - D_i(K_i - 1) - 1}{S_i}\rfloor + 1\)
        \item Par: \(p = (K_W \cdot K_H \cdot C + {\color{H3} 1}) \cdot m \ ({\color{H3} \hat{=} \ \text{Bias}})\).
        % \item Receptive field increases with depth, stride and dilation.
    \end{itemize*}
\end{definition}

\begin{definition}[Backprop]
    Kernel \(K\), input \(X\), output \(Y\), loss \(E\). Note: \(\star\) (\(A \star B \neq B \star A\)). \\
    % ROW 1: CONVOLUTION (Forward *)
    % Since forward pass flipped K, the gradient effectively "unflips" it.
    % This results in swapping the operands compared to the Correlation case.
    \(X \ast K\): 
    \(\frac{\partial E}{\partial K} = \frac{\partial E}{\partial Y} \star X = r_{180}(X \star \frac{\partial E}{\partial Y})\), 
    \(\frac{\partial E}{\partial X} = \frac{\partial E}{\partial Y} \star K\)

    % ROW 2: CORRELATION (Forward star)
    % Standard DL formulas. Input X is correlated with Error dE/dY.
    \(X \star K\): 
    \(\frac{\partial E}{\partial K} = X \star \frac{\partial E}{\partial Y}\), 
    \(\frac{\partial E}{\partial X} = \frac{\partial E}{\partial Y} \ast K\)
    \textit{Useful identity}: \(r_{180}(X) * \frac{\partial E}{\partial Y} = \frac{\partial E}{\partial Y} \star X\).
\end{definition}

\begin{definition}[Max-Pool]
    \(z_{i, j}\l l = \max\{Z_{is:is+k, js:js+k}\l{l-1}\}\) (\(s\) stride), \(\frac{\partial z_{i, j}\l l}{\partial z_{i', j'} \l{l-1}} = 1\) iff \((i',j') = (i^*, j^*)\) (max indices) and 0 otw. \(\delta\l{l-1} = \{\delta\l l\}_{i^*, j^*}\).
\end{definition}

\begin{definition}[Upsampling]
    Nearest Neighbor, Bed of nails (fill upper left with value, rest 0), Max Unpooling (remember max location), Learnable Upsampling (transpose conv: filter moves stride in output for input e, sum overlap).
\end{definition}

